\documentclass[11pt,professionalfont,aspectratio=43]{beamer}

% Assets
\usepackage[czech]{babel}				% Jazyk
\usepackage[a-2u]{pdfx}				    % Kopírování z pdfka
\usepackage{tikz}						% Schémata automatů
\usetikzlibrary{shadows, shadows.blur, shapes.geometric, arrows, positioning, overlay-beamer-styles, patterns, shadings, fadings, shapes.callouts}
% \usetikzlibrary{arrows}
\usepackage{pgfplots}
% \pgfplotsset{compat=1.15}
% \usepackage{mathrsfs}
\usepackage[export]{adjustbox}

\usepackage[T1]{fontenc}
\usefonttheme[onlymath]{serif}
% \usepackage{newtxmath}
\usepackage{pxfonts}
\usepackage{slantsc}
\usepackage{varwidth}                   % Boxy přizpůsobené šířce obsahu

\usepackage[linesnumbered,ruled,vlined]{algorithm2e}                % Pseudokód algoritmů
\usepackage[utf8]{inputenc}
% \usepackage{textcomp}
\usepackage{hyperref}
\usepackage{fancybox}

%%% FONTS
% \usefonttheme{serif}
% \usepackage{lmodern}
\usepackage{helvet}

\hypersetup{%
    pdfencoding=auto,
    pdfauthor={\insertauthor},
    pdftitle={\insertsubtitle}
}
\usepackage{float}
\usepackage{csquotes}					% české uvozovky

% Citace podle ČSN ISO 690 (biblatex-iso690). Databáze zdrojů je sdílená pro
% všechny prezentace; sazbu snímku se zdroji obstarává makro \zdrojeframe.
\usepackage[backend=biber,style=iso-numeric,sorting=none,maxnames=10]{biblatex}
\addbibresource{../assets/literatura.bib}

\usepackage{enumerate}					% enumerate environment
\usepackage{indentfirst}
\usepackage{tcolorbox}                  % Fancy boxíky
\usepackage{mathtools}
\usepackage{pifont}
\usepackage{cancel}
\usepackage{subcaption}
\usepackage{soul}
\usepackage{makecell}                   % Zalomení v buňce tabulky
\usepackage{xcolor}
\usepackage{graphicx}
\usepackage{tabularx}
\usepackage{amsmath}
\usepackage{amsthm}
\usepackage{colortbl}                   % Barvení buněk v tabulce
\usepackage{multicol}                   % Rozdělení obsahu do sloucpů

\usepackage{listings}                   % Úryvky z kódu
\definecolor{maroon}{rgb}{0.5,0,0}
\definecolor{forestgreen}{rgb}{0.13, 0.55, 0.13}

% \lstdefinelanguage{XML}
% {
%   basicstyle=\ttfamily,
%   morestring=[s]{"}{"},
%   morecomment=[s]{?}{?},
%   morecomment=[s]{!--}{--},
%   commentstyle=\color{forestgreen},
%   moredelim=[s][\color{black}]{>}{<},
%   moredelim=[s][\color{red}]{\ }{=},
%   stringstyle=\color{blue},
%   identifierstyle=\color{maroon}
% }
% Colors

\def\maincolR{0} % červená složka (0–255)
\def\maincolG{83} % zelená složka
\def\maincolB{120} % modrá složka
\definecolor{maincolor}{HTML}{1E5B4B}

% \pgfmathtruncatemacro{\authorfooterbgR}{1*\maincolR}
% \pgfmathtruncatemacro{\authorfooterbgG}{1*\maincolG}
% \pgfmathtruncatemacro{\authorfooterbgB}{1*\maincolB}
% \definecolor{authorfooterbg}{RGB}{\authorfooterbgR,\authorfooterbgG,\authorfooterbgB}

% \pgfmathtruncatemacro{\titlefooterbgR}{0.722*\maincolR}   
% \pgfmathtruncatemacro{\titlefooterbgG}{0.974*\maincolG}   
% \pgfmathtruncatemacro{\titlefooterbgB}{0.739*\maincolB}
% \definecolor{titlefooterbg}{RGB}{\titlefooterbgR,\titlefooterbgG,\titlefooterbgB}

% Beamer theme
\colorlet{themecolor}{maincolor}
\colorlet{authorfooterbg}{maincolor}
\definecolor{titlefooterbg}{HTML}{7CA096}
\definecolor{titlefooterfg}{HTML}{000000}
\definecolor{titlepagebg}{HTML}{FFFFFF}   % podklad titulní strany
\definecolor{datefooterbg}{HTML}{E3E3E3}
\colorlet{titleframeboxcolor}{maincolor}

% Background gradient
\definecolor{gradMiddle}{RGB}{212,248,255}
\definecolor{gradEnd}{RGB}{114,231,255}

% Block theme
\definecolor{blocktitlebg}{HTML}{F0DFA9}
\definecolor{blocktitlefg}{HTML}{957300}
\definecolor{blockbodybg}{HTML}{F5F0DF}

\definecolor{alertblockbodybg}{HTML}{FFE4E4}
\definecolor{alertblocktitlebg}{HTML}{ECA4A4}
\definecolor{alertblocktitlefg}{HTML}{B11313}

\definecolor{exampleblockbodybg}{HTML}{DAF0DA}
\definecolor{exampleblocktitlebg}{HTML}{BBE0BB}
\definecolor{exampleblocktitlefg}{HTML}{257A25}

\definecolor{mathfg}{HTML}{00185E}

% Názvy problémů v textu (makro \problem a odvozená v macros.tex)
\definecolor{problemfg}{HTML}{1C8C1C}

% Vývojové diagramy (makra \flowstep, \flowdecision, ... v macros.tex)
\definecolor{flowbg}{HTML}{FCE0A6}       % výplň kroků a podmínek
\definecolor{flowborder}{HTML}{A87A14}   % jejich obrys
\definecolor{flowbranchfg}{HTML}{DA0707} % popisky větví (ano/ne)

% Barevné části formulí s bublinami (styly formulapart a formulacallout
% v macros.tex): pozadí ...bg, obrys bubliny ...fg
\definecolor{calloutgreenbg}{HTML}{D2FFDC}
\definecolor{calloutgreenfg}{HTML}{078222}
\definecolor{calloutredbg}{HTML}{FFD6D2}
\definecolor{calloutredfg}{HTML}{A71000}
\definecolor{calloutbluebg}{HTML}{D6E4FF}
\definecolor{calloutbluefg}{HTML}{0A3C9E}
\definecolor{calloutorangebg}{HTML}{FFE8C6}
\definecolor{calloutorangefg}{HTML}{A65A00}
\definecolor{calloutpurplebg}{HTML}{ECDDFF}
\definecolor{calloutpurplefg}{HTML}{5B1FA7}
\definecolor{callouttealbg}{HTML}{D3ECE6}
\definecolor{callouttealfg}{HTML}{1E5B4B}

\colorlet{bulletbg}{maincolor}
% \definecolor{bulletbg}{HTML}{}

% Beamer theme
\usetheme{Madrid}
\usecolortheme{seahorse}
% \useoutertheme[subsection=false]{miniframes}

\usecolortheme[named=themecolor]{structure}
\setbeamertemplate{frame numbering}[fraction]
\setbeamertemplate{navigation symbols}{}
\setbeamercovered{invisible}

% Header/footer
\setbeamercolor{author in head/foot}{bg=authorfooterbg,fg=white}
\setbeamercolor{title in head/foot}{bg=titlefooterbg,fg=titlefooterfg}
\setbeamercolor{date in head/foot}{bg=datefooterbg,fg=black}
% Názvy problémů v zápatí přebírají barvu textu zápatí (zelená by na jeho
% podkladu zanikla)
\addtobeamertemplate{footline}{\colorlet{problemfg}{titlefooterfg}}{}
% \setbeamercolor{title}{bg=datefooterbg}
\setbeamerfont{title}{series=\bfseries}
\setbeamercolor{frametitle}{bg=datefooterbg}
% \setbeamercolor{section in head/foot}{bg=red,fg=cyan}
% \setbeamercolor{subsection in head/foot}{bg=blue,fg=yellow}

% Table of contents
\setbeamercolor{section in toc}{fg=black}
\setbeamercolor{subsection in toc}{fg=red}
\setbeamercolor{section number projected}{fg=black}

% Blocks
% Ploché bloky bez zaoblení a stínu -- definice, věty, tvrzení i příklady
% tak splývají se sazbou snímku a nepůsobí jako vystouplé "bubliny".
\setbeamertemplate{blocks}[default]

\setbeamercolor{block title}{bg=blocktitlebg, fg=blocktitlefg}
\setbeamercolor{block body}{parent=normal text,use=block title,bg=blockbodybg}

\setbeamercolor{block body alerted}{bg=alertblockbodybg}
\setbeamercolor{block title alerted}{bg=alertblocktitlebg,fg=alertblocktitlefg}

\setbeamercolor{block body example}{bg=exampleblockbodybg}
\setbeamercolor{block title example}{bg=exampleblocktitlebg,fg=exampleblocktitlefg}

% Math
\setbeamercolor{math text}{fg=mathfg}

% Itemize a enumerate
\setbeamertemplate{itemize items}[circle]
\setbeamertemplate{enumerate items}[circle]
\setbeamercolor{itemize item}{fg=bulletbg,bg=white}
\setbeamercolor{itemize subitem}{fg=bulletbg,bg=white}
\setbeamercolor{itemize subsubitem}{fg=bulletbg,bg=white}

\setbeamertemplate{enumerate items}[circle]
\setbeamercolor{item projected}{bg=bulletbg}

% U snímků s allowframebreaks se dovětek přidává až od druhého snímku, aby se
% u jednostránkového seznamu zdrojů nezobrazovalo "Zdroje I".
\setbeamertemplate{frametitle continuation}{%
    \ifnum\insertcontinuationcount>1\relax\ (pokračování)\fi
}

% Bibliografie
% Místo ikony Beameru se u položky seznamu zdrojů sází její pořadové číslo.
\setbeamertemplate{bibliography item}[text]
\setbeamercolor{bibliography item}{fg=black}
\setbeamercolor{bibliography entry author}{fg=black}
\setbeamercolor{bibliography entry title}{fg=black}
\setbeamercolor{bibliography entry location}{fg=black}
\setbeamercolor{bibliography entry note}{fg=black}

% Titulní strana je definovaná v souboru titlepage.tex.

% Enumerate
%\setlist[enumerate]{topsep=0pt,itemsep=-1ex,partopsep=1ex,parsep=1ex,label=(\arabic*)}

\MakeOuterQuote{"}

%%% Makra pro definice, věty, tvrzení, příklady, ... (vyžaduje baliček amsthm)

\theoremstyle{plain}
\newtheorem{thm}{Věta}[section]
\newtheorem{lem}[theorem]{Lemma}
\newtheorem{prop}[theorem]{Tvrzení}
\newtheorem{cor}[theorem]{Důsledek}
\newtheorem*{prop*}{Tvrzení}

\theoremstyle{definition}
\newtheorem{define}[theorem]{Definice}
% \newtheorem{example}[theorem]{Příklad}
% \newtheorem{remark}[theorem]{Poznámka}
% \newtheorem{convention}[theorem]{Úmluva}
% \newtheorem{denoting}[theorem]{Značení}

\newenvironment{defblock}[1][]{
  \begin{block}{\ifthenelse{\equal{#1}{}}
  {Definice}
  {Definice (#1)}}
}{\end{block}}
\newenvironment{thmblock}[1][]{
  \begin{block}{\ifthenelse{\equal{#1}{}}
  {Věta}
  {Věta (#1)}}
}{\end{block}}
\newenvironment{corblock}[1][]{
  \begin{block}{\ifthenelse{\equal{#1}{}}
  {Důsledek}
  {Důsledek (#1)}}
}{\end{block}}
\newenvironment{propblock}[1][]{
  \begin{block}{\ifthenelse{\equal{#1}{}}
  {Tvrzení}
  {Tvrzení (#1)}}
}{\end{block}}
\newenvironment{exblock}[1][]{
  \begin{exampleblock}{\ifthenelse{\equal{#1}{}}
  {Příklad}
  {Příklad (#1)}}
}{\end{exampleblock}}


% Colors
\definecolor{lightblue}{HTML}{009AD4}
\definecolor{lightgreen}{HTML}{68FF00}
\definecolor{darkgreen}{HTML}{00B600}
\definecolor{darkblue}{HTML}{0265b3}
\definecolor{darkred}{HTML}{DA0707}
\definecolor{lightred}{HTML}{FF5100}
\definecolor{orange}{HTML}{FFE000}
\definecolor{codeblue}{HTML}{007eed}
\definecolor{codegreen}{rgb}{0,0.6,0}
\definecolor{codegray}{rgb}{0.5,0.5,0.5}
\definecolor{codebeige}{HTML}{D4A000}
\definecolor{codepink}{HTML}{FF0055}
\definecolor{backcolour}{rgb}{0.95,0.95,0.92}

\definecolor{decproblemtitlefg}{HTML}{957300}
\definecolor{decproblembg}{HTML}{F5F0DF}

\definecolor{timportantboxcolor}{HTML}{A97100}
\definecolor{talertboxcolor}{HTML}{B70000}
\definecolor{tnoticeboxcolor}{HTML}{159C00}
\definecolor{tinfoboxcolor}{HTML}{004DFF}

\newcommand{\markred}[1]{\textcolor{lightred}{#1}}
\newcommand{\markdarkred}[1]{\textcolor{darkred}{#1}}
\newcommand{\markgreen}[1]{\textcolor{lightgreen}{#1}}
\newcommand{\markdarkgreen}[1]{\textcolor{darkgreen}{#1}}
\newcommand{\markorange}[1]{\textcolor{orange}{#1}}
\newcommand{\markblue}[1]{\textcolor{lightblue}{#1}}
\newcommand{\markdarkblue}[1]{\textcolor{darkblue}{#1}}

% Math mode settings
% \AtBeginDocument{
%     \setlength{\abovedisplayskip}{3pt}
%     \setlength{\belowdisplayskip}{-4pt}
% }

% Inline images
\newcommand{\inlineimgscale}{1.1}

% X and check mark
\newcommand{\cmark}{\markgreen{\ding{51}}}
\newcommand{\xmark}{\markred{\ding{55}}}

% Math
\newcommand{\R}{\mathbb{R}}
\newcommand{\C}{\mathbb{C}}
\newcommand{\N}{\mathbb{N}}
\newcommand{\Z}{\mathbb{Z}}
\newcommand{\Q}{\mathbb{Q}}

\newcommand{\T}[1]{#1^\top}                     % Tranpozice matice
\newcommand{\compl}[1]{#1^\complement}
\newcommand{\mapping}[3]{#1:#2 \rightarrow #3}
\newcommand{\mapto}[3]{#1:#2 \mapsto #3}
\newcommand{\set}[1]{\left\{#1\right\}}
\newcommand{\powset}{\mathcal{P}}
\newcommand{\code}[1]{\langle#1\rangle}
\DeclareMathOperator*{\argmin}{arg\,min}
\DeclareMathOperator*{\argmax}{arg\,max}

% Statistika a AI
\newcommand{\average}[1]{\overline{#1}}
\DeclareMathOperator{\median}{med}
\DeclareMathOperator{\Var}{var}
\DeclareMathOperator{\Cov}{cov}
\DeclareMathOperator{\MSE}{MSE}
\DeclareMathOperator{\SSE}{SSE}

\DeclareMathOperator{\TP}{TP}
\DeclareMathOperator{\TN}{TN}
\DeclareMathOperator{\FP}{FP}
\DeclareMathOperator{\FN}{FN}
\DeclareMathOperator{\Accuracy}{Acc}
\DeclareMathOperator{\Precision}{Prec}
\DeclareMathOperator{\Recall}{Rec}
\DeclareMathOperator{\Fscore}{F}
\DeclareMathOperator{\Gini}{Gini}

\DeclareMathOperator{\mspan}{span}
\DeclareMathOperator{\mrank}{rank}
\DeclareMathOperator{\tg}{tg}
\DeclareMathOperator{\id}{id}
\DeclareMathOperator{\dom}{dom}
\DeclareMathOperator{\rng}{rng}
\renewcommand{\emptyset}{\varnothing}
\newcommand{\binary}[1]{\ensuremath{\left\langle #1\right\rangle}_B}

% DFS in and out lists
\DeclareMathOperator{\inlist}{in}
\DeclareMathOperator{\outlist}{out}
% \DeclareMathOperator{\lowlist}{low}
\DeclareMathOperator{\key}{key}

% A* macro
\newcommand{\Astar}{$\textcolor{black}{\textup{A}^{*}}$}

% Complexity classes
\DeclareMathOperator{\classP}{P}
\DeclareMathOperator{\classNP}{NP}
\DeclareMathOperator{\classTime}{TIME}
\DeclareMathOperator{\classNTime}{NTIME}
\DeclareMathOperator{\classSpace}{SPACE}
\DeclareMathOperator{\classNSpace}{NSPACE}

%%% Názvy problémů
%
% Název problému se sází kapitálkami. Mimo matematický mód je zelený (barva
% problemfg z theme.tex), v matematickém módu se sází kapitálkami písma
% matematiky v barvě matematiky a s mezerami jako u operátoru (\mathop).
% Makra problémů fungují v obou módech:
%   Problém \SAT{} neumíme řešit rychle.     $\INDSET(G,k)=\CLIQUE(\bar{G},k)$
% Za makrem následovaným mezerou je nutné psát {} (\SAT{} je ...).
% Číslice a jiné konstanty v názvu se píší v matematickém módu ($3$-SAT).
%
%   \problem{Název}             název problému; varianta podle módu
%   \textproblem{Název}         vždy textová varianta (zelené kapitálky)
%   \mathproblem{Název}         vždy matematická varianta
%   \mathsc{Text}               kapitálky v matematickém módu (obdoba \mathrm)
%   \newproblem{\MAKRO}{Název}  zavede makro pro nový problém

\newcommand{\mathsc}[1]{\textnormal{\rmfamily\scshape #1}}
\newcommand{\textproblem}[1]{\textcolor{problemfg}{\textsc{#1}}}
\newcommand{\mathproblem}[1]{\mathop{\mathsc{#1}}\nolimits}
\DeclareRobustCommand{\problem}[1]{%
  \ifmmode\mathproblem{#1}\else\textproblem{#1}\fi}
\newcommand{\newproblem}[2]{\newcommand{#1}{\problem{#2}}}

% V záložkách a metadatech PDF zůstane jen holý název
\pdfstringdefDisableCommands{\def\problem#1{#1}}

\newproblem{\SAT}{SAT}
\newproblem{\THREESAT}{\texorpdfstring{$3$}{3}-SAT}
\newproblem{\UNSAT}{UNSAT}
\newproblem{\INDSET}{NzMna}                 % nezávislá množina
\newproblem{\CLIQUE}{Klika}
\newproblem{\SCHEDULING}{Rozvrh}
\newproblem{\SUDOKU}{Sudoku}
\newproblem{\HALT}{Halt}
\newproblem{\FACTOR}{Factor}

\DeclareMathOperator{\regex}{RegE}

%%% Ikony (vykřičník, otazník, žárovka)
%
% Lesklá koule s bílým znakem, kreslená v TikZ (vektorově, bez souborů
% s obrázky). Volitelný argument je průměr ikony:
%   \alerticon        \alerticon[2cm]        \questionicon[1.5em]     \ideaicon
% Uvnitř obrázku je jednotka rovna poloměru kruhu.

\definecolor{iconRlight}{HTML}{FFB3B3}
\definecolor{iconRmid}{HTML}{E52222}
\definecolor{iconRdark}{HTML}{B00000}
\definecolor{iconRrim}{HTML}{7A0000}
\definecolor{iconBlight}{HTML}{F0FAFF}
\definecolor{iconBmid}{HTML}{8FCBF5}
\definecolor{iconBdark}{HTML}{3D8FD6}
\definecolor{iconBrim}{HTML}{1F5C99}
\definecolor{iconGlight}{HTML}{DDFBCC}
\definecolor{iconGmid}{HTML}{52C42E}
\definecolor{iconGdark}{HTML}{218C00}
\definecolor{iconGrim}{HTML}{115500}

% Radiální stínování se středem posunutým doleva nahoru
\pgfdeclareradialshading{iconRball}{\pgfpoint{-10bp}{12bp}}{%
  color(0bp)=(iconRlight); color(14bp)=(iconRmid);
  color(40bp)=(iconRdark); color(62bp)=(iconRrim); color(100bp)=(iconRrim)}
\pgfdeclareradialshading{iconBball}{\pgfpoint{-10bp}{12bp}}{%
  color(0bp)=(iconBlight); color(14bp)=(iconBmid);
  color(40bp)=(iconBdark); color(62bp)=(iconBrim); color(100bp)=(iconBrim)}
\pgfdeclareradialshading{iconGball}{\pgfpoint{-10bp}{12bp}}{%
  color(0bp)=(iconGlight); color(14bp)=(iconGmid);
  color(40bp)=(iconGdark); color(62bp)=(iconGrim); color(100bp)=(iconGrim)}

\newlength{\iconradius}

% Pozadí: #1 = R (červené) / B (modré) / G (zelené)
\newcommand{\iconball}[1]{%
  \shade[shading=icon#1ball] (0,0) circle[radius=1];
  \draw[icon#1rim, line width=0.07\iconradius] (0,0) circle[radius=0.965];
  \fill[white, opacity=0.55, path fading=south]
    (0,0.50) ellipse[x radius=0.62, y radius=0.36];}

% Vykřičník: zužující se čárka se zaoblenými konci a tečka
\def\iconexclpath{%
  (-0.165,0.53) arc[start angle=180, end angle=0, radius=0.165]
  -- (0.075,-0.20) arc[start angle=0, end angle=-180, radius=0.075] -- cycle
  (0,-0.50) circle[radius=0.125]}
% Otazník: oblouk s krčkem se kreslí tahem, tečka zvlášť
\def\iconquesstroke{%
  (-0.33,0.31) arc[start angle=165, end angle=-28, radius=0.34]
  .. controls +(-0.06,-0.11) and (0,-0.04) .. (0,-0.17)}
\def\iconquesdot{(0,-0.58) circle[radius=0.125]}
% Žárovka: baňka s krčkem (střed baňky (0,0.22)), patice, paprsky
\def\iconbulbglass{%
  (-0.269,-0.049) arc[start angle=225, end angle=-45, radius=0.38]
  .. controls (0.20,-0.12) and (0.15,-0.20) .. (0.15,-0.28)
  -- (-0.15,-0.28) .. controls (-0.15,-0.20) and (-0.20,-0.12) .. cycle}
\def\iconbulbbase{%
  [rounded corners=0.03\iconradius] (-0.16,-0.42) rectangle (0.16,-0.33)
  (-0.16,-0.55) rectangle (0.16,-0.46)
  [sharp corners] (-0.08,-0.59) arc[start angle=180, end angle=360, radius=0.08] -- cycle}
\def\iconbulbrays{%
  \foreach \a in {0,45,90,135,180}{(0,0.22)+(\a:0.52) -- +(\a:0.66)}}

% #1 = průměr, #2 = kresba
\newcommand{\iconpicture}[2]{%
  \setlength{\iconradius}{\dimexpr #1/2\relax}%
  \begin{tikzpicture}[x=\iconradius, y=\iconradius, baseline=-0.5ex]#2\end{tikzpicture}}

\newcommand{\alerticon}[1][1cm]{\iconpicture{#1}{%
  \iconball{R}%
  \draw[iconRrim, line width=0.05\iconradius, line join=round] \iconexclpath;
  \fill[white] \iconexclpath;}}

\newcommand{\questionicon}[1][1cm]{\iconpicture{#1}{%
  \iconball{B}%
  \draw[iconBrim, line width=0.33\iconradius, line cap=round] \iconquesstroke;
  \draw[iconBrim, line width=0.05\iconradius] \iconquesdot;
  \draw[white, line width=0.23\iconradius, line cap=round] \iconquesstroke;
  \fill[white] \iconquesdot;}}

\newcommand{\ideaicon}[1][1cm]{\iconpicture{#1}{%
  \iconball{G}%
  \draw[iconGrim, line width=0.135\iconradius, line cap=round] \iconbulbrays;
  \draw[white, line width=0.08\iconradius, line cap=round] \iconbulbrays;
  \draw[iconGrim, line width=0.05\iconradius, line join=round] \iconbulbglass;
  \draw[iconGrim, line width=0.05\iconradius, line join=round] \iconbulbbase;
  \fill[white] \iconbulbglass;
  \fill[white] \iconbulbbase;}}

%%% Barevné boxíky (notice / info / important / alert)
%
% Šířka boxu se nezadává ručně -- box se přizpůsobí šířce svého obsahu.
% Obsah se vysází do prostředí varwidth, takže krátký text zůstane na
% jednom řádku; delší text se zalomí, nejvýše však na šířku \linewidth
% (tj. \textwidth mimo sloupce a minipage). Pokud je titulek širší než
% obsah, řídí šířku boxu titulek.
%
% Box může mít vlevo ikonu (\tboxicon, nastavuje prostředí; prázdné =
% bez ikony). Ikona je svisle vystředěná vůči obsahu a její šířka i mezera
% za ní se při výpočtu šířky boxu přičtou k obsahu.

\newsavebox{\tboxcontentbox}
\newsavebox{\tboxtitlebox}
\newlength{\tboxwidth}
\newlength{\tboxpadding}
\newlength{\tboxiconsize}
\newlength{\tboxiconspace}
\newlength{\tboxicongap}
\setlength{\tboxicongap}{3mm}
\def\tboxicon{}

% Geometrie boxu (musí odpovídat stylu tautobox níže):
% 2 * (boxrule + boxsep + left)
\setlength{\tboxpadding}{\dimexpr 0.5mm+1mm+4mm\relax}
\setlength{\tboxpadding}{2\tboxpadding}

\tcbset{tautobox/.style={boxrule=0.5mm, boxsep=1mm, left=4mm, right=4mm,
                         fonttitle=\bfseries}}

% Začátek sběru obsahu do boxu (velikost ikony se odvozuje od písma okolí)
\newcommand{\tboxcollect}{%
  \ifx\tboxicon\empty
    \setlength{\tboxiconspace}{0pt}%
  \else
    \setlength{\tboxiconsize}{2.6em}%
    \setlength{\tboxiconspace}{\dimexpr\tboxiconsize+\tboxicongap\relax}%
  \fi
  \begin{lrbox}{\tboxcontentbox}%
    \begin{varwidth}{\dimexpr\linewidth-\tboxpadding-\tboxiconspace\relax}%
}

% Obsah boxu včetně případné ikony
\newcommand{\tboxbody}{%
  \ifx\tboxicon\empty
    \usebox{\tboxcontentbox}%
  \else
    \hbox{%
      $\vcenter{\hbox{\tboxicon[\tboxiconsize]}}$%
      \hspace{\tboxicongap}%
      $\vcenter{\hbox{\usebox{\tboxcontentbox}}}$%
    }%
  \fi
}

% Ukončení sběru a výpočet výsledné šířky
% #1 = titulek (pro měření; prázdný u variant bez titulku)
\newcommand{\tboxmeasure}[1]{%
    \end{varwidth}%
  \end{lrbox}%
  \sbox{\tboxtitlebox}{\bfseries #1}%
  \setlength{\tboxwidth}{\dimexpr\wd\tboxcontentbox+\tboxiconspace\relax}%
  \ifdim\wd\tboxtitlebox>\tboxwidth
    \setlength{\tboxwidth}{\wd\tboxtitlebox}%
  \fi
  \addtolength{\tboxwidth}{\tboxpadding}%
  \ifdim\tboxwidth>\linewidth
    \setlength{\tboxwidth}{\linewidth}%
  \fi
}

% Vysázení boxu s titulkem: #1 = barva, #2 = titulek
\newcommand{\tboxrenderbox}[2]{%
  \tboxmeasure{#2}%
  \begin{center}%
    \begin{tcolorbox}[tautobox, colback=#1!15, colframe=#1,
                      title={#2}, width=\tboxwidth]
      \tboxbody
    \end{tcolorbox}%
  \end{center}%
}

% Vysázení boxu bez titulku: #1 = barva
\newcommand{\tboxrenderframe}[1]{%
  \tboxmeasure{}%
  \begin{center}%
    \begin{tcolorbox}[tautobox, colback=#1!15, colframe=#1,
                      width=\tboxwidth]
      \tboxbody
    \end{tcolorbox}%
  \end{center}%
}

% Poznámky (notice) -- se žárovkou pro podstatný nápad či myšlenku
\newenvironment{tnoticebox}[1]{%
  \def\tboxtitletext{#1}\def\tboxicon{\ideaicon}\tboxcollect
}{%
  \tboxrenderbox{tnoticeboxcolor}{\tboxtitletext}%
}

\newenvironment{tnoticeframe}{%
  \def\tboxicon{\ideaicon}\tboxcollect
}{%
  \tboxrenderframe{tnoticeboxcolor}%
}

% Varianty bez ikony -- ostatní poznámky
\newenvironment{tnoticebox*}[1]{%
  \def\tboxtitletext{#1}\tboxcollect
}{%
  \tboxrenderbox{tnoticeboxcolor}{\tboxtitletext}%
}

\newenvironment{tnoticeframe*}{%
  \tboxcollect
}{%
  \tboxrenderframe{tnoticeboxcolor}%
}

% Informační bloky (info) -- s otazníkem pro otázku
\newenvironment{tinfobox}[1]{%
  \def\tboxtitletext{#1}\def\tboxicon{\questionicon}\tboxcollect
}{%
  \tboxrenderbox{tinfoboxcolor}{\tboxtitletext}%
}

\newenvironment{tinfoframe}{%
  \def\tboxicon{\questionicon}\tboxcollect
}{%
  \tboxrenderframe{tinfoboxcolor}%
}

% Varianty bez ikony -- ostatní informace
\newenvironment{tinfobox*}[1]{%
  \def\tboxtitletext{#1}\tboxcollect
}{%
  \tboxrenderbox{tinfoboxcolor}{\tboxtitletext}%
}

\newenvironment{tinfoframe*}{%
  \tboxcollect
}{%
  \tboxrenderframe{tinfoboxcolor}%
}

% Důležité bloky (important)
\newenvironment{timportantbox}[1]{%
  \def\tboxtitletext{#1}\tboxcollect
}{%
  \tboxrenderbox{timportantboxcolor}{\tboxtitletext}%
}

\newenvironment{timportantframe}{%
  \tboxcollect
}{%
  \tboxrenderframe{timportantboxcolor}%
}

% Varovné bloky (alert) -- s vykřičníkem
\newenvironment{talertbox}[1]{%
  \def\tboxtitletext{#1}\def\tboxicon{\alerticon}\tboxcollect
}{%
  \tboxrenderbox{talertboxcolor}{\tboxtitletext}%
}

\newenvironment{talertframe}{%
  \def\tboxicon{\alerticon}\tboxcollect
}{%
  \tboxrenderframe{talertboxcolor}%
}

% Varianty bez ikony -- červený rámeček pro sdělení, které není přímo
% upozorněním (zdůraznění, otázka, otevřený problém, ...)
\newenvironment{talertbox*}[1]{%
  \def\tboxtitletext{#1}\tboxcollect
}{%
  \tboxrenderbox{talertboxcolor}{\tboxtitletext}%
}

\newenvironment{talertframe*}{%
  \tboxcollect
}{%
  \tboxrenderframe{talertboxcolor}%
}

% Decision problem
\newcommand{\decproblem}[3]{
    \begin{center}
        \begin{minipage}{.9\textwidth}
            \begin{table}[!htbp]
                \centering
                \bgroup
                \def\arraystretch{1.8}
                \begin{tabularx}{\textwidth}{|rX|}
                \hline
                \rowcolor{decproblembg} \multicolumn{2}{|c|}{\textcolor{decproblemtitlefg}{\textsc{#1}}} \\ \hline
                \markdarkred{\textbf{Vstup:}} & #2                \\ \hline
                \rowcolor{decproblembg} \markdarkred{\textbf{Otázka:}} & #3               \\ \hline
                \end{tabularx}
                \egroup
            \end{table}
        \end{minipage}
    \end{center}
}
\newcommand{\bigO}{\mathcal{O}}

% TODO
\newcommand{\todo}[1]{\textcolor{red}{(\noindent TODO: #1.)}}


\newcommand{\hlcode}[1]{\colorbox{red}{#1}}

% Algorithms
\numberwithin{algocf}{section}
\renewcommand{\algorithmcfname}{Algoritmus}
\SetKwInput{KwIn}{Vstup}
\SetKwInput{KwOut}{Výstup}
\SetKwProg{Function}{Funkce}{}{end}
\setlength{\algomargin}{25pt}
\renewcommand{\thealgocf}{}     % Vypnutí číslování algoritmů

% Definice stylů pro uzly
\tikzstyle{startstop} = [rectangle, rounded corners, minimum width=1.5cm, minimum height=0.5cm,text centered, draw=black, fill=red!30]
\tikzstyle{process}   = [rectangle, minimum width=1.5cm, minimum height=0.5cm, text centered, draw=black, fill=orange!30]
\tikzstyle{decision}  = [diamond, aspect=2, text centered, draw=black, fill=green!30]
\tikzstyle{arrow}     = [thick,->,>=stealth]

\tikzstyle{vertex}    = [thick, draw, circle, minimum size=4mm, inner sep=1pt, outer sep=1pt, align=center]
\tikzstyle{vertexplain}    = [minimum size=4mm, inner sep=1pt, outer sep=1pt, align=center]
\tikzstyle{verteximg}    = [draw, circle, minimum size=4mm, inner sep=0pt, outer sep=1pt, align=center, clip]
\tikzstyle{verteximgplain}    = [minimum size=4mm, inner sep=0pt, outer sep=1pt, align=center]
\tikzstyle{task}    = [draw, rounded corners=2pt, minimum width=14mm, minimum height=8mm, inner sep=2pt, align=center]

\tikzstyle{edge}    = [thick,>=stealth]
\tikzstyle{diredge}    = [thick, ->, >=stealth]

%% Automaty
\tikzstyle{state}=[%
    draw,
    circle,
    thick,
    minimum size=10mm,
    inner sep=0pt,
    outer sep=0pt
]

\tikzstyle{acceptstate}=[%
    state,
    double,
    double distance=1pt
]

\tikzstyle{transition}=[%
    ->,
    >=stealth,
    thick
]

\newcommand<>{\state}[4][]{%
    \node#5[state,#1] (#2) at #3 {#4};
}

\newcommand<>{\acceptingstate}[4][]{%
    \node#5[acceptstate,#1] (#2) at #3 {#4};
}

\newcommand<>{\transitionarrow}[5][]{%
    \path#6 (#2) edge[transition,#1] node[#3] {#4} (#5);
}

\newcommand<>{\initialstate}[6][]{%
    \node#7[state,#1] (#2) at #3 {#4};
    \draw#7[transition] #5 -- (#2.#6);
}

%%%%%

% Frames
\newcommand{\tableofcontentsframe}{
    \begin{frame}{Obsah}
        \tableofcontents
    \end{frame}
}

% \definecolor{boxcolor}{HTML}{00820a}
\newcommand{\titleframe}[1]{%
  \begin{frame}[plain]
    \begin{tikzpicture}[remember picture, overlay]
      \node at (current page.center)
        [draw=titleframeboxcolor, line width=1pt,
         inner xsep=10mm, inner ysep=4mm, rounded corners=2mm,
         text width=0.8\textwidth, align=center,
         execute at begin node=\setlength{\baselineskip}{2.2em}]
        {\Huge\markred{#1}\normalfont};
    \end{tikzpicture}
  \end{frame}
  \addtocounter{framenumber}{-1}%
}

\newenvironment{titleframeenv}[1]{
    \begin{frame}[plain]
        \begin{tcolorbox}[colback=white,boxrule=1pt]
            \begin{center}
                \Huge\markred{#1}\normalfont
            \end{center}
        \end{tcolorbox}
}{\addtocounter{framenumber}{-1}\end{frame}}

\newcommand{\icon}[1]{
  \includegraphics[height=\fontcharht\font`\B]{#1}
}

% Obrázek, který se vždy vejde na snímek. Adjustbox jej v případě potřeby
% zmenší tak, aby nepřesáhl šířku textu ani zadanou část výšky snímku; díky
% tomu sedí sazba ve všech používaných poměrech stran (4:3, 16:9 i 16:10).
% Volitelným argumentem je maximální výška (výchozí 0.68\textheight), povinným
% cesta k souboru s obrázkem (TikZ zdroj vkládaný přes \input).
\newcommand{\obrazek}[2][0.68\textheight]{%
    \begin{figure}
        \centering
        \begin{adjustbox}{max width=\linewidth, max totalheight=#1, center}
            \input{#2}
        \end{adjustbox}
    \end{figure}%
}

% Totéž pro rastrové či vektorové obrázky vkládané přes \includegraphics.
\newcommand{\obrazekgrf}[2][0.68\textheight]{%
    \begin{figure}
        \centering
        \includegraphics[width=\linewidth, max totalheight=#1, center]{#2}
    \end{figure}%
}

% Závěrečný snímek se seznamem zdrojů. Prvním (volitelným) argumentem je název
% sekce i snímku, druhým čárkou oddělený seznam klíčů z ../assets/literatura.bib.
% Citace se sázejí podle ČSN ISO 690, v pořadí, v jakém jsou klíče uvedeny.
\newcommand{\bibframe}[2][Zdroje]{%
    \section{#1}
    \nocite{#2}%
    \begin{frame}[allowframebreaks]{#1}
        \printbibliography[heading=none]
    \end{frame}
}
\lstdefinelanguage{Rust}{
    keywords={as, break, const, continue, crate, else, enum, extern, false, fn, for, if, impl, in, let, loop, match, mod, move, mut, pub, ref, return, self, Self, static, struct, super, trait, true, type, unsafe, use, where, while, async, await, dyn},
    sensitive=true,
    morecomment=[l]{//},
    morecomment=[s]{/*}{*/},
    morestring=[b]",
    morestring=[b]',
}

% Styl pro Rust
\lstdefinestyle{rust}{
    language=Rust,
    basicstyle=\footnotesize\ttfamily\color{black},
    keywordstyle=\color{codeblue}\bfseries,
    commentstyle=\color{codegreen}\ttfamily,
    stringstyle=\color{codebeige}\ttfamily,
    breaklines=true,
    xleftmargin=10pt,
    framexleftmargin=17pt,
    framexrightmargin=5pt,
    framexbottommargin=4pt,
    showspaces=false,
    showstringspaces=false,
    showtabs=false,
    tabsize=2,
}

% Styl pro Python
\lstdefinestyle{python}{
    language=Python,
    tabsize=2,
    extendedchars=true,
    breaklines=true,
    basicstyle=\footnotesize\ttfamily,
    stringstyle=\color{codebeige}\ttfamily,
    commentstyle=\color{codegreen}\ttfamily,
    keywordstyle=\color{codeblue}\bfseries,
    identifierstyle=\color{black},
    showspaces=false,
    showtabs=true,
    xleftmargin=10pt,
    framexleftmargin=17pt,
    framexrightmargin=17pt,
    framexbottommargin=4pt,
    showstringspaces=false,
    morekeywords={False, None, True, and, as, assert, async, await, break, class, continue, def, del, elif, else, except, finally, for, from, global, if, import, in, is, lambda, nonlocal, not, or, pass, raise, return, try, while, with, yield}
}

% Styl pro C#
\lstdefinestyle{sharpc}{
    language={csh},
    tabsize=2,
    extendedchars=true,
    breaklines=true,
    basicstyle=\footnotesize\ttfamily,
    stringstyle=\color{codebeige}\ttfamily,
    commentstyle=\color{codegreen}\ttfamily,
    keywordstyle=\color{codeblue}\bfseries,
    identifierstyle=\color{black},
    showspaces=false,
    showtabs=true,
    xleftmargin=10pt,
    framexleftmargin=17pt,
    framexrightmargin=17pt,
    framexbottommargin=4pt,
    showstringspaces=false,
    morekeywords={abstract, event, new, struct, as, explicit, null, switch, base, extern, object, this, bool, false, operator, throw, break, finally, out, true, byte, fixed, override, try, case, float, params, typeof, catch, for, private, uint, char, foreach, protected, ulong, checked, goto, public, unchecked, class, if, readonly, unsafe, const, implicit, ref, ushort, continue, in, return, using, decimal, int, sbyte, virtual, default, interface, sealed, volatile, delegate, internal, short, void, do, is, sizeof, while, double, lock, stackalloc, else, long, static, enum, namespace, string}
}

% Styl pro C
\lstdefinestyle{clang}{
    language=C,
    basicstyle=\scriptsize\ttfamily\color{black},
    keywordstyle=\color{codeblue}\bfseries,
    commentstyle=\color{codegreen}\ttfamily,
    numberstyle=\tiny\color{codegray},
    stringstyle=\color{codebeige}\ttfamily,
    breakatwhitespace=false,
    breaklines=true,
    captionpos=b,
    keepspaces=true,
    numbersep=5pt,
    showspaces=false,
    showstringspaces=false,
    showtabs=false,
    morekeywords={
        auto, break, case, char, const, continue, default, do, double, else, enum, extern, float, for, goto, if, inline, int, long, register, restrict, return, short, signed, sizeof, static, struct, switch, typedef, union, unsigned, void, volatile, while, 
        _Alignas, _Alignof, _Atomic, _Bool, _Complex, _Generic, _Imaginary, _Noreturn, _Static_assert, _Thread_local
    },
    tabsize=2,
    moredelim=*[l][\color{red}\bfseries]{\#},
    identifierstyle=\color{black},
    xleftmargin=10pt,          % přidá odsazení zleva
    framexleftmargin=17pt,     % odsazení rámce zleva, pokud používáte rámeček
    framexrightmargin=5pt,     % odsazení rámce zprava
    framexbottommargin=4pt     % odsazení rámce zdola
}

% Nastavení výchozího stylu
\lstset{
    escapeinside={(*}{*)},
    style=clang,
    basicstyle=\footnotesize\ttfamily\color{black}
}

\def\presentationtitle{Problém 3-SAT, problém nezávislé množiny a~kliky}
\def\presentationtitlefooter{3-SAT, NzMna a~klika}

\begin{document}

    % =====================================================================
% Titulní strana prezentací
%
% Stranu dělí šikmá hrana: horní barevná plocha nese název předmětu
% a téma přednášky, spodní světlá část autora, kontakt, školu, školní
% rok a logo školy.
%
% Šablona není vázaná na konkrétní předmět -- všechny údaje se berou
% z příkazů \title, \subtitle, \author, \email, \institute, \date
% a z makra \schoollogo.
% =====================================================================

\makeatletter
\@ifundefined{presentationtitlefooter}{%
  \let\presentationtitlefooter\presentationtitle
}{}
\makeatother

% Název předmětu lze v konkrétní prezentaci přepsat pomocí
% \def\subjectname{...} před vložením tohoto souboru.
\providecommand{\subjectname}{Teoretická informatika}
% Cesta k logu školy (relativně ke složce prezentace).
\providecommand{\schoollogo}{../images/SPSE-Jecna_Logo.pdf}

\title[\presentationtitlefooter]{\subjectname}
\subtitle{\presentationtitle}
\author{David Weber}
\def\email{weber3@spsejecna.cz}
\institute{SPŠE Ječná}
\date{%
    Rok
    \pgfmathsetmacro{\firstyear}{ifthenelse(\month<9, \year-1, \year)}%
    \pgfmathsetmacro{\secondyear}{ifthenelse(\month<9, \year, \year+1)}%
    \pgfmathtruncatemacro{\secondyearlasttwo}{mod(\secondyear,100)}%
    \pgfmathprintnumber[fixed,precision=0,zerofill,1000 sep={}]{\firstyear}/%
    \pgfmathprintnumber[fixed,precision=0,zerofill,1000 sep={}]{\secondyearlasttwo}%
}

% --- Pomocná makra ---------------------------------------------------

% Vypnutí dělení slov -- na titulní straně působí rušivě. Kromě
% penalizace je potřeba i \emergencystretch, jinak TeX slovo rozdělí
% ve chvíli, kdy jinou přijatelnou sazbu nenajde.
\newcommand{\tpnohyphen}{%
    \hyphenpenalty=10000\exhyphenpenalty=10000%
    \tolerance=9999\emergencystretch=3em\relax
}

% Téma přednášky. Pozor: \tpnohyphen musí být mimo skupinu s fontem --
% parametry odstavce se čtou až při jeho ukončení. Nastavení
% \hyphenchar na -1 dělení vypne i tam, kde by TeX jinak neměl jiné
% řešení.
\newcommand{\tptopic}[1]{%
    \tpnohyphen
    {\usebeamerfont{title}\Large
     \hyphenchar\font=-1\relax
     #1\par
     \hyphenchar\font=`\-\relax}%
}

% Text sázený na barevném podkladu: matematika, názvy problémů i makro
% \Astar musí být bílé, ne tmavě modré, zelené, resp. černé.
\makeatletter
\newcommand{\tpinverse}{%
    \tpnohyphen
    \setbeamercolor{math text}{fg=white}%
    \colorlet{problemfg}{white}%
    \@ifundefined{Astar}{}{\renewcommand{\Astar}{\textup{A}\textsuperscript{*}}}%
}
\makeatother

% Školní rok. Čísla se sázejí v matematickém režimu, takže barvu je
% nutné nastavit i pro "math text" -- jinak zůstanou tmavě modrá.
% #1 = barva
\newcommand{\tpdate}[1]{{\setbeamercolor{math text}{fg=#1}\color{#1}\insertdate}}

% --- Titulní strana --------------------------------------------------

% Vnější obrázek slouží pouze k umístění kresby do levého horního rohu
% strany. Vlastní kresba je vnořený obrázek s vlastní soustavou
% souřadnic: (0,0) je levý horní roh strany, (1,-1) pravý dolní.
% Odkazy na "current page" uvnitř cest by ve zdejší sestavě balíčků
% vedly k chybám "Missing number", proto se používají jen v kotvě uzlu.
\begin{frame}[plain]
    \begin{tikzpicture}[remember picture, overlay]
      \node[anchor=north west, inner sep=0pt] at (current page.north west) {%
        \begin{tikzpicture}[x=\paperwidth, y=\paperheight]
            % Podklad
            \fill[titlepagebg] (0,0) rectangle (1,-1);

            % Světlejší doprovodný pruh podél šikmé hrany
            \fill[titlefooterbg]
                (0,-0.615) -- (1,-0.435) -- (1,-0.475) -- (0,-0.655) -- cycle;

            % Barevná plocha se šikmou hranou
            \fill[maincolor]
                (0,0) -- (1,0) -- (1,-0.42) -- (0,-0.60) -- cycle;

            % Název předmětu
            \node[anchor=north west, white, font=\normalsize\bfseries\scshape,
                  text width=0.80\paperwidth, align=left, inner sep=0pt]
                at (0.09,-0.12) {\tpinverse\inserttitle};

            % Téma přednášky
            \node[anchor=north west, white, text width=0.70\paperwidth,
                  align=left, inner sep=0pt]
                at (0.09,-0.19) {\tpinverse\tptopic{\insertsubtitle}};

            % Autor, kontakt, škola
            \node[anchor=south west, text width=0.55\paperwidth,
                  align=left, inner sep=0pt]
                at (0.09,-0.84) {%
                    {\usebeamerfont{author}\insertauthor\par}%
                    {\small\texttt{\email}\par}%
                    \vskip0.5em
                    {\small\insertinstitute\par}%
                };

            % Školní rok
            \node[anchor=south west, font=\small, inner sep=0pt]
                at (0.09,-0.91) {\tpdate{black}};

            % Logo školy
            \node[anchor=south east, inner sep=0pt]
                at (0.91,-0.90)
                {\includegraphics[width=0.14\paperwidth]{\schoollogo}};
        \end{tikzpicture}};
    \end{tikzpicture}%
\end{frame}


    % =================================================================
    \section{Problém 3-SAT}
    % =================================================================

    \begin{titleframeenv}{Problém $\THREESAT$}
    \end{titleframeenv}

    \begin{frame}{Problém $\THREESAT$}
        \begin{itemize}
            \visible<2->{\item Problém $\SAT$ neumíme řešit v~polynomiálním čase.}
            \visible<3->{\item Zkusíme proto původní problém trochu zjednodušit, a~to tak, že budeme uvažovat formule v~CNF, kde každá klauzule má maximálně 3 literály.}
        \end{itemize}
        \visible<4->{\begin{defblock}[3-CNF]
            Formule $\varphi$ je v~3-CNF právě tehdy, když je v~CNF a~každá klauzule obsahuje maximálně 3 literály.
        \end{defblock}}
        \visible<5->{\decproblem{Problém $\THREESAT$}{
            Formule $\varphi$ v~3-CNF.
        }{
            Je $\varphi$ splnitelná?
        }}
    \end{frame}

    \begin{frame}{Problém $\THREESAT$}
        \visible<2->{\begin{tinfoframe*}
            \begin{center}
                Ukážeme, že problém $\THREESAT$ je "stejně těžký" jako původní $\SAT$.
            \end{center}
        \end{tinfoframe*}}
        \begin{itemize}
            \visible<3->{\item Jeden problém jsme schopni "převést" na druhý.}
            \visible<4->{\item Přesněji k~zadané formuli $\varphi$ v~CNF jsme schopni najít formuli $\psi$ v~3-CNF, která je splnitelná právě tehdy, když je i~původní formule $\varphi$ splnitelná.}
        \end{itemize}
        \visible<5->{\begin{timportantframe}
            \begin{center}
                Převod mezi formulemi však musíme umět provést \textbf{rychle}!
            \end{center}
        \end{timportantframe}}
    \end{frame}

    % =================================================================
    \section{Polynomiální převoditelnost problémů}
    % =================================================================

    \begin{titleframeenv}{Převoditelnost problémů}
    \end{titleframeenv}

    \subsection{Zavedení polynomiální převoditelnosti}

    \begin{frame}{Polynomiální převoditelnost}
        \begin{itemize}
            \visible<2->{\item Formálně lze problém chápat jako zobrazení binárního kódování vstupu do množiny $\set{0,1}$.}
            \visible<3->{\item Existuje celá řada problémů, avšak lze mezi nimi najít zajímavé vztahy.}
            \visible<4->{\item Resp. některé problémy nemusíme nutně umět řešit přímo $\implies$ místo toho je převedeme na ty, které řešit umíme!}
        \end{itemize}
        \visible<5->{\begin{defblock}[Polynomiální převoditelnost]
            Nechť $A,B$ jsou rozhodovací problémy. Říkáme, že $A$ lze \markdarkred{převést} na $B$, píšeme $A\rightarrow B$, právě tehdy, když existuje zobrazení $\mapping{f}{\set{0,1}^{*}}{\set{0,1}^{*}}$ takové, že
            \[\forall x\in\set{0,1}^{*}: A(x)=B(f(x)),\]
            a~navíc $f(x)$ lze spočítat v~polynomiálním čase. Funkci $f$ říkáme \markdarkred{převod} nebo také \markdarkred{redukce}.
        \end{defblock}}
    \end{frame}

    \subsection{Význam}

    \begin{frame}{Co to znamená?}
        \begin{itemize}
            \visible<2->{\item Problém $A$ je převoditelný na problém $B$, pokud každý vstup $x$ pro $A$ jsme schopni \textbf{rychle sestrojit} odpovídající vstup $f(x)$ pro $B$.}
            \visible<3->{\item Přitom musí platit, že odpověď $A$ na $x$ musí být stejná jako odpověď $B$ na $f(x)$.}
            \visible<4->{\item V~našem případě: Máme $\varphi$ v~CNF a~chceme v~polynomiálím čase zkonstruovat formuli $\psi$ v~3-CNF takovou, aby
            \[\SAT(\varphi)=\THREESAT(\psi).\]}
            \visible<5->{\item Tato konstrukce představuje ono zobrazení $f$, tj. $f(\varphi)=\psi$, nebo-li
            \[\SAT(\varphi)=\THREESAT(f(\varphi)).\]}
        \end{itemize}
    \end{frame}

    \subsection{Poznámky k~převoditelnosti}

    \begin{frame}{Poznámky k~převoditelnosti}
        \begin{itemize}
            \visible<2->{\item Nemusí nutně platit, že problémy $A$ a~$B$ jsou na sebe \textbf{vzájemně} převoditelné.}
            \visible<3->{\begin{center}
                Tj. $A\rightarrow B$ nutně ještě neznamená, že i~$B\rightarrow A$.
            \end{center}}
            \visible<4->{\item Např. uvažujme dvojici jednoduchých problémů $K,L$ definovaných následovně:}
            \begin{align*}
                \visible<5->{K(x)&=0\qquad\text{Na každý vstup odpověz nulou},\\}
                \visible<6->{L(x)&=1\qquad\text{Na každý vstup odpověz jedničkou}.}
            \end{align*}
            \visible<7->{\begin{talertframe*}
                \begin{center}
                    Takovou dvojici problémů na sebe nepřevedeme \textbf{ani jedním směrem}.
                \end{center}
            \end{talertframe*}}
        \end{itemize}
    \end{frame}

    \subsection{Příklad: Převod mezi SAT a~HALT}

    \begin{frame}{Převod $\SAT\rightarrow\HALT$}
        \visible<2->{Převod $\SAT$ na $\HALT$ lze provést jednoduše: Sestrojíme program $P$ se vstupem $\varphi$, který}
        \begin{enumerate}
            \visible<3->{\item postupně zkouší všechna ohodnocení proměnných formule $\varphi$,}
            \visible<4->{\item pokud najde splňující ohodnocení, \textbf{zastaví se},}
            \visible<5->{\item jinak vstoupí do nekonečného cyklu.}
        \end{enumerate}
        \visible<6->{Text programu jsme schopni sestrojit \textbf{v~lineárním čase} vzhledem k~délce formule $\varphi$.}

        \visible<7->{\begin{tnoticeframe*}
            Tedy skutečně platí, že
            \[\varphi\;\text{je splnitelná} \iff \text{program $P$ se nad $\varphi$ zastaví},\]
            nebo-li $\SAT(\varphi)=\HALT(P,\varphi)$.
        \end{tnoticeframe*}}
    \end{frame}

    \begin{frame}{Převod $\HALT\rightarrow\SAT$ není možný}
        \begin{itemize}
            \visible<2->{\item Už jsme si zdůvodnili, že ne pro všechny páry $(P,x)$ lze rozhodnout, zda se program $P$ nad $x$ zastaví.}
            \visible<3->{\item Konkrétně myšlenka byla založena na konstrukci programu $D$, který invertoval výstup $\HALT$.}
            \visible<4->{\item Pro spor předpokládejme, že umíme převést \textbf{problém zastavení} na \textbf{problém splnitelnosti formule v~CNF}.}
            \begin{enumerate}
                \visible<5->{\item Instanci $(P,x)$ převedeme polynomiálně na formuli $\varphi=f(P,x)$.}
                \visible<6->{\item Následně rozhodneme, zda $\SAT(f(P,x))=1$, tj. zdali je formule $f(P,x)$ splnitelná.}
            \end{enumerate}
        \end{itemize}
        \visible<7->{\begin{talertframe*}
            \begin{center}
                Jenže pro $P=D$ nelze problém zastavení rozhodnout, zatímco v~tomto případě bychom jednoduše mohli sestrojit formuli $\psi=f(D,D)$ a~rozhodnout, zdali $\SAT(\psi)=1$. \markdarkred{To je spor!}
            \end{center}
        \end{talertframe*}}
    \end{frame}

    % =================================================================
    \section{Převod mezi SAT a~3-SAT}
    % =================================================================

    \begin{titleframeenv}{Převod mezi $\SAT$ a~$\THREESAT$}
    \end{titleframeenv}

    \subsection{Převod 3-SAT na SAT}

    \begin{frame}{Převod $\THREESAT\rightarrow\SAT$}
        \begin{itemize}
            \visible<2->{\item Na začátku máme formuli $\varphi$ v~3-CNF a~chceme k~ní sestrojit formuli $\psi$ v~CNF, která je splnitelná právě tehdy, když je splnitelná i~$\varphi$.}
            \visible<3->{\begin{talertframe*}
                \begin{center}
                    Každá formule ve 3-CNF je zároveň i~v CNF $\implies$ ponecháme původní formuli $\varphi$.
                \end{center}
            \end{talertframe*}}
            \visible<4->{\item Za převodní zobrazení $f$ stačí zvolit identitu, tj. $f(x)=x$.}
        \end{itemize}
    \end{frame}

    \subsection{Převod SAT na 3-SAT}

    \begin{frame}{Převod $\SAT\rightarrow\THREESAT$}
        \begin{itemize}
            \visible<2->{\item Zde je již zjevné, že ne každá formule $\varphi$ v~CNF je automaticky i~v 3-CNF.}
            \visible<3->{\item U~vstupní formule proto musíme rozdělit "dlouhé" klauzule na menší.}
            \visible<4->{\item Na vstupu tedy máme formuli $\varphi$ v~CNF. Mohou nastat 2 scénáře:}
            \begin{enumerate}
                \visible<5->{\item Formule $\varphi$ neobsahuje klauzuli s~více než 3 literály. Pak jsme hotovi. \cmark}
                \visible<6->{\item Jinak vybereme nějakou klauzuli $C$, která obsahuje $k>3$ literálů, a~rozdělíme ji na dvě části $\alpha,\beta$:
                \[C=(\underbrace{\ell_1\lor \ell_2}_\alpha\lor\overbrace{\ell_3\lor\ldots\lor\ell_k}^\beta).\]
                Tedy $\alpha=\ell_1\lor\ell_2$ a~$\beta=\ell_3\lor\ldots\lor\ell_k$.}
            \end{enumerate}
        \end{itemize}
    \end{frame}

    \begin{frame}{Nové klauzule s~$\alpha$ a~$\beta$}
        \begin{itemize}
            \visible<2->{\item Přidáme novou proměnnou $x$ a~původní klauzuli zapíšeme pomocí dvojice nových klauzulí následovně:
            \[(\ell_1\lor \ell_2\lor\ldots\lor\ell_k)\equiv(\alpha\lor x)\land(\beta\lor\neg x).\]}
            \visible<3->{\item Všimněme si, že}
            \begin{itemize}
                \visible<4->{\item Klauzule $(\alpha\lor x)=(\ell_1\lor \ell_2\lor x)$ má 3 literály.}
                \visible<5->{\item Klauzule $(\beta\lor\neg x)=(\ell_3\lor\ldots\lor\ell_k\lor\neg x)$ má $k-1$ literálů.}
            \end{itemize}
            \visible<6->{\begin{tnoticeframe}
                \begin{center}
                    Původní klauzuli jsme rozdělili na \textbf{dvojici kratších klauzulí}.
                \end{center}
            \end{tnoticeframe}}
            \visible<7->{\item Pokud je klauzule $(\beta\lor\neg x)$ stále dlouhá, tak postup opakujeme.}
        \end{itemize}
    \end{frame}

    \begin{frame}{Hodnota proměnné $x$}
        \begin{itemize}
            \item Nyní zbývá dořešit, jak nastavit hodnotu proměnné $x$.
            \begin{enumerate}
                \item Pokud $\alpha=1$, pak je klauzule $(\alpha\lor x)$ splněna $\implies$ nastavením $x=0$ splníme i~druhou klauzuli $(\beta\lor\neg x)$.
                \item Pokud $\alpha=0$, pak nutně $\beta=1$. Jinak by původní klauzule
                \[(\ell_1\lor\ldots\lor\ell_k)=(\alpha\lor\beta)\]
                nebyla splněna. $\implies$ Nastavíme $x=1$, abychom splnili klauzuli $(\alpha\lor x)$.
            \end{enumerate}
        \end{itemize}
    \end{frame}

    \begin{frame}{Závěrem k~převodu $\SAT\rightarrow\THREESAT$}
        \begin{itemize}
            \visible<2->{\item Opakovanou aplikací tohoto postupu dostaneme formuli $\psi$, která je v~3-CNF a~je splnitelná právě tehdy, když i~$\varphi$ je splnitelná.}
            \visible<3->{\item Rozdělení klauzule
            \[C=(\ell_1\lor\ell_2\lor \ldots\lor \ell_k)\]
            délky $k$ provedeme v~$k-3$ krocích. Tím postupně dostaneme
            \begin{align*}
                C&\equiv(\ell_1\lor\ell_2\lor x_1)\land(\neg x_1\lor\ell_3\lor x_2)\land\dots\land(\neg x_{k-4}\lor\ell_{k-2}\lor x_{k-3})\\ 
                &\land(\neg x_{k-3}\lor\ell_{k-1}\lor\ell_k),
            \end{align*}
            kde $x_1,x_2,\dots,x_{k-3}$ jsou nové proměnné.}
            \visible<4->{\item Tento rozklad potrvá lineárně dlouho vzhledem k~délce klauzule, tj. $\bigO(k)$.}
            \visible<5->{\item Celkově převod všech dlouhých klauzulí tedy potrvá $\bigO(n)$, kde $n$ je počet literálů.}
        \end{itemize}
    \end{frame}

    \begin{frame}{$\SAT$ a~$\THREESAT$ jsou vzájemně převoditelné}
        \obrazek{images/sat-3sat.tex}
    \end{frame}

    \begin{frame}{Převod $\SAT\rightarrow\THREESAT$}{Příklad}
        \small
        Mějme formuli $\varphi$ v~CNF s~jednou klauzulí:
        \[\varphi=a\lor b\lor \neg c\lor \neg b\lor d\lor \neg d.\]
        \begin{itemize}
            \item Rozdělíme $\varphi$ na dvě klauzule pomocí nové proměnné $x$:
            \[(a\lor b\lor \markdarkred{x})\land(\neg c\lor \neg b\lor d\lor \neg d\lor \markdarkred{\neg x}).\]
            \item Pro druhou klauzuli nyní opakujeme stejný postup: Zavedeme si novou proměnnou $y$ a~rozdělíme ji na dvojici nových klauzulí, tj.
            \[(a\lor b\lor x)\land(\neg c\lor \neg b\lor \markdarkred{y})\land(d\lor \neg d\lor \neg x\lor \markdarkred{\neg y}).\]
        \end{itemize}
    \end{frame}

    \begin{frame}{Převod $\SAT\rightarrow\THREESAT$}{Příklad}
        \small
        \begin{itemize}
            \item A~nyní ještě poslední klauzuli rozdělíme pomocí nové proměnné $z$:
            \[(a\lor b\lor x)\land(\neg c\lor \neg b\lor y)\land(d\lor \neg d\lor \markdarkred{z})\land(\neg x\lor \neg y\lor \markdarkred{\neg z}).\]
            \item Všechny klauzule už mají nejvýše 3~literály, takže jsme hotovi. \cmark
        \end{itemize}
    \end{frame}

    \begin{frame}{Převod $\SAT\rightarrow\THREESAT$}{Příklad}
        \visible<2->{Výsledná formule v~3-CNF je tedy
        \[\psi=(a\lor b\lor x)\land(\neg c\lor \neg b\lor y)\land(d\lor \neg d\lor z)\land(\neg x\lor \neg y\lor \neg z).\]}
        \visible<3->{Zkusíme ohodnocení proměnných $a=0,b=1,c=1,d=1$. Jak musíme nastavit $x,y,z$?}
        \begin{itemize}
            \visible<4->{\item První klauzule: $a\lor b = 0\lor 1=1$, tedy nastavíme $x=0$.}
            \visible<5->{\item Druhá klauzule: $\neg c\lor \neg b=\neg 1\lor \neg 1=0$, nastavíme $y=1$.}
            \visible<6->{\item Třetí klauzule: $d\lor \neg d=1\lor\neg 1=1$, nastavíme $z=0$.}
            \visible<7->{\item Čtvrtá klauzule: $(\neg x\lor \neg y\lor \neg z)=\neg 0\lor \neg 1\lor\neg 0=1$.}
        \end{itemize}
        \visible<8->{Platí, že $\varphi(\ldots)=\psi(\ldots)=1$.}
    \end{frame}

    % =================================================================
    \section{Obtížnost problémů}
    % =================================================================

    \begin{titleframeenv}{Obtížnost problémů}
    \end{titleframeenv}

    \begin{frame}{Převoditelnost a~"obtížnost" problému}
        \begin{itemize}
            \visible<2->{\item $A\rightarrow B$ v~podstatě říká, že problém $B$ je alespoň tak těžký jako problém $A$.}
            \visible<3->{\begin{tnoticeframe}
                \begin{center}
                    Pokud umíme vyřešit $B$, pak vyřešit $A$ je nanejvýše "polynomiálně obtížnější".
                \end{center}
            \end{tnoticeframe}}
            \visible<4->{\item Zároveň to znamená, že pomocí problému $B$ můžeme vyřešit $A$ s~maximálně polynomiálním zpomalením (kvůli převodu vstupu).}
            \visible<5->{\item Mnemotechnická pomůcka: \emph{Obtížnost teče ve směru šipky.}}
        \end{itemize}
    \end{frame}

    \begin{frame}{Umíme-li $B$ rychle, umíme i~$A$ rychle}
        \visible<2->{\begin{propblock}
            Pokud $A\rightarrow B$ a~$B$ lze řešit v~polynomiálním čase, pak i~$A$ lze řešit v~polynomiálním čase.
        \end{propblock}}
        \begin{itemize}
            \visible<3->{\item Uvažujme algoritmus řešící problém $B$ v~čase $\bigO(b^{k})$, kde $b$ je délka vstupu tohoto problému a~$k\in\R^{+}$.}
            \visible<4->{\item Dále mějme zobrazení $f$ převádějící $A$ na $B$ v~čase $\bigO(a^{\ell})$ pro vstup délky $a$ a~konstantu $\ell\in\R^{+}$.}
            \visible<5->{\item Chceme-li spočítat $A(x)$ pro vstup $x$ délky $a$, spočítáme nejprve $f(x)$ v~čase $\bigO(a^{\ell})$}
            \visible<6->{\begin{center}
                $\implies$ Tím dostaneme výstup délky $\bigO(a^{\ell})$ (delší bychom ani nestihli vypsat).
            \end{center}}
            \visible<7->{\item Tento výstup $f(x)$ předáme na vstup algoritmu řešící problém $B$, který nad ním stráví $\bigO((a^{\ell})^{k})=\bigO(a^{k\ell}).$}
            \visible<8->{\item Celkově strávíme výpočtem $\bigO(a^{\ell}+a^{k\ell})$, což je polynom v~délce původního vstupu $x$.}
        \end{itemize}
    \end{frame}

    \begin{frame}{Znázornění tvrzení}
        \obrazek{images/obtiznost.tex}
    \end{frame}

    \begin{frame}{Co z~toho plyne?}
        \visible<2->{\begin{tnoticeframe*}
            \begin{center}
                Speciálně to znamená, že problém $\THREESAT$ a~$\SAT$ jsou \textbf{stejně těžké problémy}.
            \end{center}
        \end{tnoticeframe*}}
        \visible<3->{\begin{tinfoframe*}
            \begin{center}
                Zkrácení délky klauzulí pro nás bude mít jednu zásadní výhodu u~dalších problémů.
            \end{center}
        \end{tinfoframe*}}
    \end{frame}

    \begin{frame}{Tranzitivita převoditelnosti}
        \visible<2->{\begin{propblock}
            Nechť $A,B,C$ jsou rozhodovací problémy, přičemž $A\rightarrow B$ a~$B\rightarrow C$. Pak platí $A\rightarrow C$.
        \end{propblock}}
        \begin{itemize}
            \visible<3->{\item Z~definice máme, že existují polynomiálně vyčíslitelná zobrazení $f,g$ taková, že
            \[A(x)=B(f(x))\quad\text{a}\quad B(x)=C(g(x))\;,\; x\in\set{0,1}^{*}.\]}
            \visible<4->{\item Z~toho tedy plyne, že $A(x)=C(g(f(x)))$.}
            \visible<5->{\item Definujeme zobrazení $h$ jako $h(x)=g(f(x))$.}
            \visible<6->{\item Jsme-li schopni spočítat $f$ v~čase $\bigO(a^{k})$ a~$g$ v~čase $\bigO(b^{\ell})$, kde $a,b$ jsou délky vstupů a~$k,\ell\in\R^{+}$, pak zobrazení $h$ spočítáme v~čase
            \[\bigO((a^{k})^{\ell})=\bigO(a^{k\ell}),\]
            což je polynom v~délce původního vstupu $x$.}
        \end{itemize}
    \end{frame}

    % =================================================================
    \section{Problém nezávislé množiny}
    % =================================================================

    \begin{titleframeenv}{Problém nezávislé množiny}
    \end{titleframeenv}

    \begin{frame}{Nezávislá množina}
        \visible<2->{\begin{defblock}[Nezávislá množina]
            Mějme libovolný graf $G=(V,E)$. Množina vrcholů $M\subseteq V$ je \markdarkred{nezávislá}, pokud žádné dva vrcholy $u,v\in M$ nejsou spojeny hranou.
        \end{defblock}}
        \visible<3->{\obrazek[0.53\textheight]{images/nzmna.tex}}
    \end{frame}

    \subsection{Existence nezávislé množiny}

    \begin{frame}{Existence nezávislé množiny}
        \visible<2->{Na samotnou existenci nezávislé množiny nemá moc smysl se ptát $\implies$ stačí zkrátka vzít \textbf{jednovrcholovou} nebo \textbf{prázdnou} množinu.}
        \visible<3->{\begin{timportantframe}
            \begin{center}
                Zajímavější je však otázka, zdali existuje nezávislá množina \textbf{určité velikosti}.
            \end{center}
        \end{timportantframe}}
    \end{frame}

    \begin{frame}{Nezávislá množina}{Příklad: Úplný graf}
        \obrazek{images/uplny-graf-nzmna.tex}
        \visible<2->{\begin{center}
            Největší nezávislá množina v~úplném grafu obsahuje pouze jeden vrchol.
        \end{center}}
    \end{frame}

    \begin{frame}{Nezávislá množina}{Příklad: Cyklus}
        \obrazek{images/cyklus-nzmna.tex}
        \visible<2->{\begin{center}
            Největší nezávislá množina v~cyklu $C_n$ má velikost $\lfloor n/2\rfloor$.
        \end{center}}
    \end{frame}

    \begin{frame}{Nezávislá množina}{Příklad: Strom}
        \obrazek{images/strom-nzmna.tex}
        \visible<2->{V~obecném stromu $T$ na $n$ vrcholech platí pro největší nezávislou množinu $N$ následující odhad:
        \[\left\lceil\frac{n}{2}\right\rceil\leqslant|N|\leqslant n-1.\]}
    \end{frame}

    \subsection{Problém nezávislé množiny v~grafu}

    \begin{frame}{Problém nezávislé množiny v~grafu}
        \visible<2->{\decproblem{Problém nezávislé množiny v~grafu ($\NZMNA$)}{
            Graf $G$ a~číslo $k\in\N$.
        }{
            Existuje v~grafu $G$ nezávislá množina velikosti alespoň $k$?
        }}
        \visible<3->{\begin{talertframe*}
            \begin{center}
                Ukážeme, že problém $\NZMNA$ je \textbf{stejně těžký} jako $\SAT$.
            \end{center}
        \end{talertframe*}}
    \end{frame}

    \subsection{Převod 3-SAT na NzMna}

    \begin{frame}{Převod $\THREESAT\rightarrow\NZMNA$}
        \begin{itemize}
            \visible<2->{\item Máme zadanou formuli $\varphi$ v~3-CNF s~$n$ literály. Chceme pro ni sestrojit graf $G$ a~číslo $k$ takové, že platí:
            \begin{center}
                $\varphi$ je splnitelná právě tehdy, když v~grafu $G$ existuje nezávislá množina velikosti alespoň $k$.
            \end{center}}
            \visible<3->{\item Označme si jako $k$ počet klauzulí formule $\varphi$.}
            \visible<4->{\item Z~každé klauzule vybereme jeden literál, který ji bude splňovat.}
            \visible<5->{\begin{talertframe}
                \begin{center}
                    \textbf{Pozor! Literály nesmíme vybírat konfliktně}. Když vybereme z~jedné klauzule literál $\ell$, nemůžeme pak z~jiné vybrat $\neg\ell$.
                \end{center}
            \end{talertframe}}
        \end{itemize}
    \end{frame}

    \begin{frame}{Převod klauzulí na trojcykly}
        \begin{itemize}
            \visible<2->{\item Každou klauzuli s~$i$ literály, kde $i\in\set{1,2,3}$, převedeme na úplný graf $K_i$. Tzn. pro každý literál $\ell$ zavedeme jeden vrchol.}
            \visible<3->{\obrazek{images/prevod-3sat-nzmna.tex}}
            \visible<4->{\item Nakonec spojíme hranou ještě všechny literály (příslušné vrcholy) s~jejich negacemi.}
            \visible<5->{\item Tím jsme sestrojili příslušný graf $G$.}
        \end{itemize}
    \end{frame}

    \begin{frame}{Převod klauzulí na trojcykly}{Příklad}
        \obrazek{images/prevod-3sat-nzmna-priklad.tex}
    \end{frame}

    \begin{frame}{Korektnost převodu}
        \begin{itemize}
            \visible<2->{\item Uvažujme, že existuje splňující ohodnocení formule $\varphi$.}
            \visible<3->{\item Vybereme z~každé klauzule \textbf{jeden literál}, který ji splňuje.}
        \end{itemize}
        \visible<4->{\begin{tnoticeframe*}
            \begin{center}
                Příslušné vrcholy vybraných literálů tvoří nezávislou množinu velikosti $k$ v~grafu $G$.
            \end{center}
        \end{tnoticeframe*}}
        \visible<5->{Polynomialita převodu je zřejmá: Graf $G$ má maximálně $n$ vrcholů a~hran může mít maximálně \textbf{kvadraticky mnoho}.}
    \end{frame}

    \subsection{Převod NzMna na SAT}

    \begin{frame}{Převod $\NZMNA\rightarrow\SAT$}{Bonus}
        \visible<2->{\begin{tinfoframe*}
            \begin{center}
                Bonusová část. Postup je převzat z~knihy \cite{mares-valla-labyrint}.
            \end{center}
        \end{tinfoframe*}}
        \begin{itemize}
            \visible<3->{\item Máme zadaný graf $G=(V,E)$ s~vrcholy očíslovanými $1,\dots,n$ a~číslo $k\in\N$. Chceme sestrojit formuli $\varphi$ v~CNF takovou, že platí:
            \begin{center}
                $\varphi$ je splnitelná právě tehdy, když v~grafu $G$ existuje nezávislá množina velikosti alespoň $k$.
            \end{center}}
            \visible<4->{\item Pro každý vrchol $j$ zavedeme proměnnou $v_j$ s~významem
            \[v_j=1\iff\text{vrchol $j$ vybereme do množiny}.\]}
        \end{itemize}
        \visible<5->{\begin{tnoticeframe*}
            \begin{center}
                Ohodnocení proměnných $v_1,\dots,v_n$ tak popisuje množinu vrcholů $M=\set{j\mid v_j=1}$.
            \end{center}
        \end{tnoticeframe*}}
    \end{frame}

    \begin{frame}{Množina musí být nezávislá}
        \begin{itemize}
            \visible<2->{\item Žádné dva vybrané vrcholy nesmí být spojeny hranou.}
            \visible<3->{\item Pro každou hranu $ij\in E$ proto přidáme klauzuli
            \[(\neg v_i\lor\neg v_j).\]}
            \visible<4->{\item Klauzule je splněna právě tehdy, když alespoň jeden z~vrcholů $i,j$ \textbf{není} vybrán.}
        \end{itemize}
        \visible<5->{\begin{timportantframe}
            \begin{center}
                Zbývá zajistit, aby byla množina \textbf{dostatečně velká}.
            \end{center}
        \end{timportantframe}}
        \visible<6->{\begin{talertframe}
            \begin{center}
                Formule však neumí přímo "spočítat", kolik proměnných $v_j$ je rovno jedné.
            \end{center}
        \end{talertframe}}
    \end{frame}

    \begin{frame}{Velikost množiny: lest s~očíslováním}
        \begin{itemize}
            \visible<2->{\item Vybrané vrcholy \textbf{očíslujeme} čísly $1,\dots,k$.}
            \visible<3->{\item Zavedeme matici proměnných $X=(x_{i,j})$ typu $k\times n$, kde
            \[x_{i,j}=1\iff\text{$i$-tým prvkem množiny je vrchol $j$}.\]}
            \visible<4->{\item Např. pro $n=4$ a~$k=2$ (řádky odpovídají pořadí prvku, sloupce vrcholům):
            \[X=\begin{pmatrix}
                0 & \markdarkred{1} & 0 & 0\\
                0 & 0 & 0 & \markdarkred{1}
            \end{pmatrix}.\]}
            \visible<5->{\item Prvním prvkem množiny je vrchol $2$, druhým vrchol $4$.}
        \end{itemize}
        \visible<6->{\begin{tnoticeframe*}
            \begin{center}
                Je-li v~každém řádku \textbf{právě jedna} jednička a~v~každém sloupci \textbf{nejvýše jedna}, vybírá matice $k$ různých vrcholů.
            \end{center}
        \end{tnoticeframe*}}
    \end{frame}

    \begin{frame}{Podmínky kladené na matici $X$}
        \small
        \setlength{\abovedisplayskip}{3pt}
        \setlength{\belowdisplayskip}{3pt}
        \begin{enumerate}
            \visible<2->{\item V~každém řádku je \textbf{nejvýše jedna} jednička:
            \[(\neg x_{i,j}\lor\neg x_{i,j'})\quad\text{pro všechna $i$ a~$j\neq j'$}.\]}
            \visible<3->{\item V~každém řádku je \textbf{alespoň jedna} jednička:
            \[(x_{i,1}\lor x_{i,2}\lor\dots\lor x_{i,n})\quad\text{pro všechna $i$}.\]}
            \visible<4->{\item V~každém sloupci je \textbf{nejvýše jedna} jednička:
            \[(\neg x_{i,j}\lor\neg x_{i',j})\quad\text{pro všechna $j$ a~$i\neq i'$}.\]}
            \visible<5->{\item Očíslované vrcholy musí být vybrané, tj. $x_{i,j}\Rightarrow v_j$. Implikaci $a\Rightarrow b$ zapíšeme jako klauzuli $(\neg a\lor b)$:
            \[(\neg x_{i,j}\lor v_j)\quad\text{pro všechna $i$ a~$j$}.\]}
        \end{enumerate}
    \end{frame}

    \begin{frame}{Výsledná formule}
        \visible<2->{Formule $\varphi$ je konjunkcí všech dosud zavedených klauzulí:}
        \visible<2->{\obrazek[0.66\textheight]{images/nzmna-sat-formule.tex}}
        \visible<3->{\begin{center}
            \footnotesize Formule $\varphi$ je v~CNF s~proměnnými $v_1,\dots,v_n$ a~$x_{1,1},\dots,x_{k,n}$.
        \end{center}}
    \end{frame}

    \begin{frame}{Korektnost převodu}
        \small
        \begin{itemize}
            \visible<2->{\item[$\Rightarrow$] Nechť $M$ je nezávislá množina, $|M|\geqslant k$. Vybereme z~ní $k$ vrcholů $u_1,\dots,u_k$.}
            \begin{itemize}
                \visible<3->{\item Nastavíme $v_j=1\iff j\in M$ a~$x_{i,j}=1\iff j=u_i$.}
                \visible<4->{\item Snadno ověříme, že jsou splněny všechny klauzule.}
            \end{itemize}
            \visible<5->{\item[$\Leftarrow$] Nechť máme splňující ohodnocení formule $\varphi$ a~$M=\set{j\mid v_j=1}$.}
            \begin{itemize}
                \visible<6->{\item Díky klauzulím pro hrany je $M$ nezávislá.}
                \visible<7->{\item Každý řádek matice $X$ určuje jeden vrchol a~různé řádky určují různé vrcholy.}
                \visible<8->{\item Všechny tyto vrcholy leží v~$M$ $\implies$ $|M|\geqslant k$.}
            \end{itemize}
        \end{itemize}
        \visible<9->{\begin{tinfoframe*}
            \begin{center}
                Množina $M$ může obsahovat i~\textbf{neočíslované} vrcholy. To nevadí, podstatné je, že obsahuje alespoň $k$ očíslovaných.
            \end{center}
        \end{tinfoframe*}}
    \end{frame}

    \begin{frame}{Převod je polynomiální}
        \begin{itemize}
            \visible<2->{\item Je-li $k>n$, nezávislá množina velikosti $k$ neexistuje $\implies$ vypíšeme nesplnitelnou formuli, např. $(y)\land(\neg y)$.}
            \visible<3->{\item Jinak platí $k\leqslant n$ a~počty klauzulí jsou:}
            \begin{itemize}
                \visible<4->{\item nezávislost: $|E|\leqslant n^{2}$,}
                \visible<5->{\item řádky: nejvýše $k\cdot n^{2}$ dvojic a~$k$ dlouhých klauzulí délky $n$,}
                \visible<6->{\item sloupce: nejvýše $n\cdot k^{2}$,}
                \visible<7->{\item vazba na $v_j$: $k\cdot n$.}
            \end{itemize}
            \visible<8->{\item Celkem tedy formule obsahuje $\bigO(n^{3})$ literálů.}
        \end{itemize}
        \visible<9->{\begin{tnoticeframe*}
            \begin{center}
                Formuli sestrojíme v~polynomiálním čase $\implies$ $\NZMNA\rightarrow\SAT$.
            \end{center}
        \end{tnoticeframe*}}
    \end{frame}

    \begin{frame}{Převod $\NZMNA\rightarrow\SAT$}{Příklad}
        \begin{columns}[c]
            \begin{column}{0.36\textwidth}
                \obrazek[0.36\textheight]{images/nzmna-sat-priklad.tex}
            \end{column}
            \begin{column}{0.6\textwidth}
                \visible<2->{Hledáme nezávislou množinu velikosti alespoň $k=2$.}

                \visible<3->{Klauzule pro hrany:
                \begin{align*}
                    &(\neg v_1\lor\neg v_2)\land(\neg v_1\lor\neg v_3)\land(\neg v_1\lor\neg v_4)\\
                    \land\;&(\neg v_2\lor\neg v_3)\land(\neg v_3\lor\neg v_4)
                \end{align*}}
            \end{column}
        \end{columns}
        \begin{itemize}
            \visible<4->{\item Zvolíme nezávislou množinu $M=\set{2,4}$, tj. $v_2=v_4=1$ a~$v_1=v_3=0$.}
            \visible<5->{\item Každá klauzule pro hranu obsahuje $\neg v_1$ nebo $\neg v_3$ $\implies$ je splněna.}
        \end{itemize}
    \end{frame}

    \begin{frame}{Převod $\NZMNA\rightarrow\SAT$}{Příklad}
        \begin{itemize}
            \visible<2->{\item Vrchol $2$ očíslujeme jako první, vrchol $4$ jako druhý prvek množiny:
            \[X=\begin{pmatrix}
                0 & \markdarkred{1} & 0 & 0\\
                0 & 0 & 0 & \markdarkred{1}
            \end{pmatrix},\quad\text{tj. } x_{1,2}=x_{2,4}=1.\]}
            \visible<3->{\item V~každém řádku je právě jedna jednička a~v~každém sloupci nejvýše jedna. \cmark}
            \visible<4->{\item Vazba na $v_j$: implikace $x_{1,2}\Rightarrow v_2$ a~$x_{2,4}\Rightarrow v_4$ platí, ostatní mají nepravdivý předpoklad. \cmark}
        \end{itemize}
        \visible<5->{\begin{tnoticeframe*}
            \begin{center}
                Našli jsme splňující ohodnocení formule $\varphi$, které odpovídá nezávislé množině $\set{2,4}$.
            \end{center}
        \end{tnoticeframe*}}
    \end{frame}

    \subsection{Převod NzMna na 3-SAT}

    \begin{frame}{Převod $\NZMNA\rightarrow\THREESAT$}
        \begin{itemize}
            \visible<2->{\item Převod dostaneme složením dvou převodů, které už známe.}
            \visible<3->{\item Problém nejprve převedeme na $\SAT$ (viz bonusová část) a~pak využijeme převod $\SAT$ na $\THREESAT$. Díky tranzitivitě převoditelnosti platí $\NZMNA\rightarrow\THREESAT$.}
        \end{itemize}
        \visible<4->{\obrazek{images/nzmna-3sat.tex}}
    \end{frame}

    % =================================================================
    \section{Problém kliky v~grafu}
    % =================================================================

    \begin{titleframeenv}{Problém kliky v~grafu}
    \end{titleframeenv}

    \subsection{Klika v~grafu}

    \begin{frame}{Klika v~grafu}
        \visible<2->{\begin{defblock}[Klika]
            Nechť $G=(V,E)$ je libovolný graf. Pak \markdarkred{klikou} nazveme libovolný úplný podgraf $H=(V^{\prime},E^{\prime})$ grafu $G$, tzn. $V^{\prime}\subseteq V$ a~$E^{\prime}\subseteq E$.
        \end{defblock}}
        \visible<3->{\obrazek[0.53\textheight]{images/klika.tex}}
    \end{frame}

    \begin{frame}{Ještě jeden příklad}
        \obrazek[0.79\textheight]{images/klika-2.tex}
    \end{frame}

    \subsection{Doplněk grafu}

    \begin{frame}{Doplněk grafu}
        \visible<2->{\begin{defblock}[Doplněk grafu]
            \markdarkred{Doplňkem} grafu $G=(V,E)$ je graf $\bar{G}=(V,E^{\prime})$, kde pro každé $u,v\in V$ platí, že
            \[uv\in E^{\prime}\iff uv\notin E,\quad u\neq v.\]
        \end{defblock}}
        \visible<3->{Tzn. doplněk grafu $G$ obsahuje všechny hrany, které nejsou obsaženy v~$G$.}
        \visible<4->{\obrazek[0.32\textheight]{images/doplnek.tex}}
    \end{frame}

    \subsection{Problém kliky v~grafu}

    \begin{frame}{Problém kliky v~grafu}
        \visible<2->{\decproblem{Problém kliky v~grafu ($\KLIKA$)}{
            Graf $G$ a~číslo $k\in\N$.
        }{
            Obsahuje $G$ kliku o~alespoň $k$ vrcholech?
        }}
        \visible<3->{\begin{talertframe*}
            \begin{center}
                Problém kliky v~grafu \textbf{je komplementární} k~problému nezávislé množiny v~grafu.
            \end{center}
        \end{talertframe*}}
    \end{frame}

    \subsection{Převody mezi Klika a~NzMna}

    \begin{frame}{Převody mezi $\KLIKA$ a~$\NZMNA$}
        \begin{itemize}
            \visible<2->{\item Jeden problém lze převést na druhý jednoduše sestrojením \textbf{doplňku grafu}.}
            \visible<3->{\item Jsou-li libovolné vrcholy spojeny v~$G$ hranou, pak v~$\bar{G}$ již nejsou spojeny hranou.}
            \visible<4->{\begin{tnoticeframe}
                \begin{center}
                    Z~kliky vznikne nezávislá množina stejné velikosti a~naopak.
                \end{center}
            \end{tnoticeframe}}
            \visible<5->{\obrazek[0.38\textheight]{images/klika-nzmna.tex}}
        \end{itemize}
    \end{frame}

    \begin{frame}{Převody mezi $\KLIKA$ a~$\NZMNA$}
        \begin{itemize}
            \visible<2->{\item Máme zadaný graf $G$ a~číslo $k\in\N$.}
            \visible<3->{\item Sestrojíme doplněk $\bar{G}$ grafu $G$ a~ponecháme číslo $k$.}
            \visible<4->{\item Pak v~grafu $G$ existuje klika velikosti alespoň $k$ právě tehdy, když v~$\bar{G}$ existuje nezávislá množina velikosti alespoň $k$.}
        \end{itemize}
        \visible<5->{\begin{tinfoframe*}
            \begin{center}
                Formálně jsme sestrojili převod $f$ definovaný předpisem $f(G,k)=(\bar{G},k)$.
            \end{center}
        \end{tinfoframe*}}
        \visible<6->{Zjevně platí
        \[\KLIKA(G,k)=\NZMNA(\bar{G},k),\;\text{resp.}\;\NZMNA(G,k)=\KLIKA(\bar{G},k).\]}
    \end{frame}

    \begin{frame}{Převod je polynomiální}
        \visible<2->{Vrcholy původního grafu ponecháváme, přičemž hran má graf maximálně $\bigO(n^{2})$, kde $n$ je počet vrcholů.}
        \visible<3->{\begin{tnoticeframe*}
            \begin{center}
                Tzn. převod "oběma směry" potrvá \textbf{polynomiálně dlouho}.
            \end{center}
        \end{tnoticeframe*}}
    \end{frame}

    % =================================================================
    \section{Shrnutí}
    % =================================================================

    \begin{frame}{Shrnutí}
        \visible<2->{\obrazek{images/shrnuti.tex}}
        \visible<3->{\begin{talertframe*}
            \begin{center}
                Převod $\NZMNA\rightarrow\SAT$ je součástí bonusové části.
            \end{center}
        \end{talertframe*}}
    \end{frame}

    \zdrojeframe{papadimitriou-complexity, mares-valla-labyrint}

\end{document}